\chapter{Вычислительная лаборатория серии}
\label{ch:computational-laboratory}

\section{Затянувшаяся лаборатория, а не только теоремы}
\label{sec:lab-role}

Монография задаёт закон~$g$, геометрию носителя и цепочку $M\to T$.
Проверка этой цепочки для читателя устроена как \emph{затянувшаяся физическая лаборатория},
в которой вместо осциллографа --- симулятор решётки и алгебраические сверки,
а вместо лабораторного журнала --- явные \textbf{записи опытов} в конце каждого тома.

Три роли не смешиваются:
\begin{description}[style=nextline,leftmargin=3.2cm]
\item[\textbf{Теория (MODEL)}]
  что \emph{должно} следовать из посылок; без дат сессий и без статуса кода.
\item[\textbf{Аппарат}]
  дискретная реализация шага~$g$, приборы осреднения, сценарии начальных условий
  (в репозитории проекта --- пакет симуляции и сценарии запуска).
\item[\textbf{Журнал прогонов}]
  серия вычислительных экспериментов: полная постановка, наблюдаемые,
  анализ по \cref{ch:measurement-processing}, вывод и статус (закрыт / открыт / провал).
\end{description}

\begin{remark}[Не подмена натурного опыта]
Вычислительный прогон проверяет \emph{согласованность} заявленного $g$ и моста $M\to T$
на заданной постановке.
Сравнение с PDG/CODATA остаётся на~$\Tlayer$ и требует метрологии эталона
(\cref{ch:measurement-processing}), а не «красивого скриншота поля''.
\end{remark}

\section{Бремя доказательства}
\label{sec:lab-burden-of-proof}

Утверждения серии сильные, потому что вывод жёсткий: \textbf{полное квантование на~$M$},
две скорости $c_0$ и $c$, лестница $E_0$, перепись $M=97$ --- не метафора и не «авторы решили, что ок''.
Публикация сама по себе никому ничего не доказывает: для каждого класса утверждений мы обязаны дать
хотя бы одно из:
\begin{enumerate}[label=\textbf{B\arabic*:},leftmargin=*]
\item \textbf{доказательство} из явных посылок (теорема с гиперссылкой на аксиому);
\item \textbf{запись опыта} (серия CE-*) с наблюдаемой, порогом PASS/FAIL и командой воспроизведения;
\item \textbf{явный статус «открыто''} с критерием decisive-провала (что именно сломает модель).
\end{enumerate}
Сравнение с CODATA/PDG без строки Verify и без \cref{ch:measurement-processing} --- недопустимая
замена опыта.
Рецензент не обязан «догадываться'', что имелось в виду: если в тексте нет CE-записи или реестра,
утверждение считается \emph{не проверенным в рамках серии}, даже если формула выглядит правдоподобно.

\section{Единый протокол записи опыта}
\label{sec:lab-protocol}

Каждая запись --- один файл \texttt{book/sources/lab-journal/<id>.tex}: обычный \LaTeX{},
четыре подзаголовка \texttt{\textbackslash subsubsection*\{Постановка\}} и т.\,д., без макросов журнала.
В \texttt{lab-*-records.tex} только \texttt{\textbackslash input\{lab-journal/ce-m-00.tex\}}.

\begin{labrecord}{Шаблон (пустая форма)}{lab:template}
\subsubsection*{Постановка}
Физический вопрос и место опыта в MODEL (слои $M$ / $T$, не реестр тестов).
Что проверяем по смыслу; рабочая гипотеза; какие величины снимаем; критерий согласия с теорией (пороги погрешности, монотонность, масштаб).

\subsubsection*{Теоретическая часть}
Определения, теоремы, следствия из книги --- зачем именно эти наблюдаемые должны вести себя так; обязательная проверка размерностей всех входящих и выходящих величин в системе СИ ($\mathrm{L, M, T}$).

\subsubsection*{Ход эксперимента}
Воспроизводимые шаги (коммит, команда прогона при необходимости), таблицы с числами.

\subsubsection*{Анализ, расчёт погрешности и выводы}
Интерпретация таблиц; тип A/B (\cref{ch:measurement-processing}); физический вывод; статус записи в журнале.

\end{labrecord}

Серии опытов (\texttt{labseries}) объединяют прогоны с общим аппаратом
(один и тот же кипящий вакуум, одна цепочка аннигиляции, одна лестница~$\alpha$).

\section{Связь с метрологией и критическими экспериментами}
\label{sec:lab-metrology-link}

\begin{enumerate}[label=\textbf{L\arabic*:},leftmargin=*]
\item \textbf{Постановка} фиксирует, что является наблюдаемой, а что --- артефактом дискретизации.
\item \textbf{Анализ} --- по \cref{ch:measurement-processing}: прямые/косвенные величины, GUM, Стьюдент для ансамбля.
\item \textbf{Критические эксперименты} в конце тома --- не буллет-лист, а развёрнутый журнал
  (\CiteConstruction{ch:critical-construction}; журнал этажей --- том \emph{floors}; аналоги в прочих томах).
\item Опыт со статусом \emph{открыт} остаётся в журнале с явным критерием decisive-провала модели.
\end{enumerate}

\section{Классы вычислительных опытов}
\label{sec:lab-classes}

\begin{description}[style=nextline,leftmargin=3.4cm]
\item[\textbf{Алгебра / замыкание}]
  SI-строка, лестница $\alpha$, тождества масс без поля на решётке.
  Критерий --- внутренняя согласованность и $E_n$ к CODATA при заявленных $u$.
\item[\textbf{Макро-поле на карте}]
  дисперсия, фронт, скорость $c$ после $\mathcal{B}$ (\CiteConstruction{ch:macro}).
\item[\textbf{Составной узел на этаже}]
  два- и $n$-телесные сценарии, оболочки, аннигиляция, прокси диполя и времени жизни.
\item[\textbf{Отрицательные / контроль}]
  сценарий, который \emph{не} должен давать эффект (контроль pp вместо pm и т.\,п.).
\end{description}

Журнал пополняется по мере закрытия UoW; закрытая запись не переписывает MODEL,
а фиксирует, какой класс проверок на данном этапе считается выполненным.

\section{Формальное операторное ядро в Lean~4 (\texttt{lean/MtCarrier/}, 97 теорем от уравнения Каталана $1+2^d=d^2$)}
\label{sec:lab-lean4-operator-kernel}

Все фундаментальные определения и теоремы модели формализованы и проверены компилятором \textbf{Lean~4} в пакете \texttt{lean/MtCarrier/} (\texttt{Spec.lean}, \texttt{FaceAlgebra.lean}, \texttt{HolonomyMaxwell.lean}, \texttt{TrajectoryIV263.lean}, \texttt{UnifiedOperators.lean}, \texttt{AxiomAudit.lean} и исполняемый мост \texttt{DumpJson.lean}, сверяемый побитово с \texttt{mt\_ca/carrier\_core.py} в проверке \texttt{Lean4\_carrier\_bridge}).

Архитектура формального ядра подчиняется четырём принципам:
\begin{enumerate}[leftmargin=*,noitemsep]
\item \textbf{Единый диофантов корень Каталана:} размерность $d\in\mathbb{N}$ остаётся свободной переменной во всех операторах и связывается исключительно операторным предикатом $\mathrm{CatalanLock}(d)\equiv (1+2^d=d^2)$ и универсальным комбинатором $\mathrm{CatalanLaw}(P,d)\equiv (\mathrm{CatalanLock}(d)\leftrightarrow \mathrm{CatalanLock}(d)\wedge P(d))$ (\texttt{catalan\_lock\_iff\_d3}, \texttt{catalan\_law\_of\_d3}).
\item \textbf{Инкапсулированный 3D-автомат (\texttt{runCarrierAutomatonOp}):} эволюция поля $\Psi^{(t)} = \widehat{\mathcal{G}}_d^{[t]}(\Psi^{(0)})$ (\texttt{iterFieldStepOp}) инкапсулирует суммирование по оболочке кубооктаэдра $O_2(d)=\mathrm{cubicOrbit3D}(d-1)$, кусочно-линейные проекции граней $\mathcal{C}_d,\mathcal{S}_d$, знаменатель импеданса $\mathrm{gateDenomOp}$, деление вниз \texttt{Int.fdiv} и мёртвую зону затвора $\mathrm{heisenbergDeadbandOp}$, предоставляя чистый интерфейс $\text{Вход}\to\text{Выход}$ (\texttt{encapsulated\_automaton\_input\_output\_theorem}).
\item \textbf{Универсальная линейность, мультилинейность и гомоморфизм внешней среды (голографический экран Маркова):} граничные операторы среды $(\mathrm{envSumMuListOp},\mathrm{envZrListOp},\mathrm{envZiListOp})$ являются точными аддитивными гомоморфизмами конкатенации многокомпонентных потоков среды $\mathcal{E}_1\mathbin{++}\mathcal{E}_2$ (\texttt{multi\_component\_environment\_homomorphism\_and\_markov\_blanket}, \texttt{universal\_operator\_linearity\_and\_multilinearity}), а при масштабном переходе на этаж $h_L(K)=K\ell_P$ квадратичный фактор $K^2$ однородно факторизуется в числителе и знаменателе ворот (\texttt{scale\_invariance\_of\_holonomy\_and\_gate\_operators}).
\item \textbf{Межэтажный изоморфизм замкнутых оболочек и полное квантование (\texttt{cross\_floor\_shell\_closure\_isomorphism\_operator}, \texttt{complete\_quantization\_of\_phase\_space\_and\_all\_observables\_operator}):} собственные оболочки кубооктаэдра $O_2(d)$ единым операторным законом порождают лестницу Белла--Цирельсона со знаменателем $(\caSignalSpeed/\speedC)^2=\mathrm{c0OverCSqOp}(d)$, баллистическую передачу энергии и информации без потерь в защищённом волноводе Гейзенберга, вихревой атлас частиц, кварков, бозонов и матриц PMNS/CKM на кликах $\mathcal{C}_{13}(d)$, периодическую систему вплоть до 8-го периода с четырьмя физическими пределами и рецептами синтеза $Z=119,120$ и $N_{\mathrm{island}}=184$, кавитированные гиперядра, топологию «ядерной пасты» нейтронных звёзд, несингулярное ядро чёрной дыры Паули ($D_{\min}=\varepsilon(d)>0$), ряд сжатия Гаусса--Грина, планетогенез всех типов планетных ядер, космобиологию трёх окон жизни $m_{\mathrm{env}}\in\{-1,0,+1\}$, предсказуемый мутагенез и сквозной обратный компилятор 3D-фенотипа в последовательность ДНК.
\end{enumerate}
